\documentclass[12pt]{article}
\usepackage[utf8]{inputenc}
\usepackage{graphicx}
\usepackage{hyperref}
\usepackage{listings}
\usepackage{float}
\usepackage{siunitx}
\usepackage[ngerman]{babel}
\usepackage{csquotes}
\usepackage[
    backend=biber,
    style=ieee,
    sorting=none
]{biblatex}
\addbibresource{literatur.bib}
\usepackage{amsmath}
\usepackage[
    top=2.5cm,
    bottom=2.5cm,
    left=3cm,
    right=3cm,
    headheight=14pt 
]{geometry}
\usepackage{setspace}
\usepackage{fancyhdr}
\usepackage[locale=DE]{siunitx}
% \usepackage{minted} for code

%  Logo in der Ecke 
%\pagestyle{fancy}
%\fancyhf{} 
%\fancyhead[R]{\includegraphics[width=3cm]{Image.jpg}} 
%\fancyfoot[C]{\thepage} 
%\renewcommand{\headrulewidth}{0pt}
%\renewcommand{\footrulewidth}{0pt} 
%\geometry{
%  top=2cm,
%  %left=2.5cm,
%  %right=2.5cm,
%  bottom=2cm,
%  %headheight=1.5cm % Passt die Höhe der Kopfzeile an, damit das Logo Platz hat
%}


\pagenumbering{gobble}
\begin{document}
\begin{titlepage}

	% Fehler "destination with the same identifier" unterdrücken...
  \setcounter{page}{-1}

	% Titelseite
	\begin{figure}[h]
		\begin{minipage}[b]{12 cm}
			{\sffamily
				{\Large HfT Stuttgart } \\
				Fakultät  C
			}
		\end{minipage}
	\end{figure}

	\vfill
	
	\begin{center}
		% Seminararbeit 
		\noindent { \huge
			Interdisziplinäres Projekt \\
			{Flexible Roboterprogrammierung für Montage und Demontage}\\
		}
		\vspace{14mm}
		% Titel
		\noindent \textbf{\huge
		  Abschlussbericht\\
		}
		\vspace{20mm}	
  \noindent { 
			Prof. Dr.-Ing. Darko Katic\\
		}
	\end{center}
	
	\vfill
	
\noindent\textbf{Valentin Bäuerle}

\noindent\rule{\textwidth}{0.4mm}

\noindent\href{mailto:21bava1bif@hft-stuttgart.de}{21bava1bif@hft-stuttgart.de}

\noindent Informatik

\noindent Matr.-Nr. 1004286

\end{titlepage}

\newpage
\section*{Eidesstattliche Erklärung}
%todo
Hiermit erkläre ich Valentin Bäuerle (1004286), dass ich den vorliegenden Abschlussbericht selbstständig und ohne fremde Hilfe bzw. unerlaubte Hilfsmittel angefertigt, andere als die angegebenen Quellen und Hilfsmittel nicht benutzt und die den benutzten Quellen wörtlich oder inhaltlich entnommenen Stellen als solche kenntlich gemacht habe.

\section*{Sperrvermerk}
Die vorliegende Arbeit beinhaltet interne und vertrauliche Informationen der HFT Stuttgart.
Valentin Bäuerle verpflichtet sich, alle im Rahmen der Arbeit gewonnenen Informationen und Erkenntnisse für einen Zeitraum von 5 Jahren vertraulich zu behandeln, d.h. sie nicht Dritten zugänglich zu machen oder anderweitig zu verwerten. 
Dies gilt auch für Familienangehörige oder Personen, die bei der Erstellung der Arbeit – sei es durch Korrekturlesen oder ähnliches - behilflich sind. Die Arbeit darf den betreuenden Professoren ausschließlich zur fachlichen Bewertung vorgelegt werden.
Die Weitergabe des Inhalts der Arbeit im Gesamten oder in Teilen sowie das Anfertigen von Kopien oder Abschriften – auch in digitaler Form – sind grundsätzlich untersagt. Publikationen jeglicher Art (gilt auch für die Aufnahme im hochschuleigenen Bibliotheksbestand) bedürfen der schriftlichen Genehmigung der HFT Stuttgart.\\ \\

%todo
\noindent Waiblingen, den 15.05.2026

\begin{figure}[h!] % Optionale Positionierungsparameter
    \includegraphics[width=0.5\textwidth]{images/Unterschrift.png} % Pfad und Dateiname deines Bildes
    \label{fig:meinbild} % Ein eindeutiges Label, um auf das Bild zu verweisen
\end{figure}

\noindent ……………………………………

\noindent Unterschrift (Valentin Bäuerle)


\newpage

\vspace*{-2cm}  % Wert anpassen bis es passt
\tableofcontents

\newpage

\pagenumbering{arabic}

\section{Einleitung}
In diesem Kapitel wird der Kontext der vorliegenden Arbeit beschrieben uund das Thema begündet. Zunächst wird die Forschungsfrage und die Ziele der Arbeit erläutert. Anschließend wird die Methodik der Arbeit vorgestellt und die Aufbau der Arbeit erläutert.
%Ausgangssituation, Forschungsfrage, Ziel der Arbeit, Methodische Vorgehensweise, Aufbau der Arbeit

%\subsection{Einstieg}
\subsection{Ausgangssituation}
Bekanntes 3-10 sätze Was gibt es bereits zum thema ?
Definitionen 1-2 sätze pro Def. 
relevante Details und aspekte 10-20 sätze. worum geht es, welche fragen, was spielt eine rolle

Das manuelle Montieren und Demontieren von Bauteilen ist ein zentrales Element moderner Fertigungsprozesse, insbesondere bei der Montage von Elektronikkomponenten, der Automobilproduktion und der Luft- und Raumfahrt \cite{manuelle_montage}. Da manuelle Montageprozesse zeit- und kostenintensiv sind, ist deren Automatisierung ein zentrales Ziel in der Robotik \cite{automation_aufwand}, um Effizienz, Präzision und Sicherheit zu erhöhen. Roboterarme werden zunehmend eingesetzt, um diese Aufgaben zu übernehmen\cite{robotics_manufacturing_replacement}.

Hierbei besteht die Herausforderung darin, dass XXXXX

Diese Einschränkungen motivieren den Einsatz XXXXXX 

Im Folgenden wird die daraus resultierende Problemstellung und Motivaiton dieser Arbeit näher erläutert.
	
\subsection{Forschungsfrage}

\subsection{Ziel der Arbeit}

\subsection{Methodische Vorgehensweise}

\subsection{Aufbau des Berichts}
Welche inhalte finden sich im ersten kapitel?
welche inhalte stehen im zweiten kapitel zum forschungsstand?
welche aspekte werden im dritten kapitel behandelt?
Welche inhalte finden sich im vierten kapitel, dem Ergebniskapitel?
Welche inhalte werden im letzen Kapitel behandelt?

Der vorliegende Bericht gliedert sich in vier Kapitel. Kapitel 2 legt die theoretischen Grundlagen dar: Es führt den Dobot Magician als Hardwareplattform ein, beschreibt die verwendete Simulationsumgebung und definiert zentrale Begriffe wie Flexibilität im Kontext roboterbasierter Montage. Kapitel 3 stellt die betrachteten Verfahren – klassische Planungsansätze, lernbasierte Methoden sowie Foundation Models – systematisch vor und bewertet sie anhand definierter Kriterien in einer Bewertungsmatrix. Kapitel 4 fasst die Ergebnisse des Vergleichs zusammen und leitet daraus eine begründete Empfehlung für die vertiefte Untersuchung im Rahmen der Bachelorarbeit ab.

\newpage

\section{Stand der Forschung und Technik}
\subsection{Robotische Montage und Demontage in der Praxis}
\subsection{Computervision und automatische Objekterkennung in der Robotik}
\subsection{Mensch-Roboter-Kollaboration und kooperative Fertigungsprozesse}
\subsection{Forschungslücke und Zielstellung dieser Arbeit}

\newpage

\section{Grundlagen}
Dieses Kapitel beschreibt die technischen und theoretischen Grundlagen der vorliegenden Arbeit. Zunächst wird der verwendete Roboterarm Dobot Magician als experimentelle Plattform vorgestellt. Anschließend werden 
%todo:die eingesetzte Simulationsumgebung sowie die Definition von Flexibilität im Kontext der Arbeit erläutert. Abschließend erfolgt eine Einordnung von Montage- und Demontageaufgaben als Robotikproblem. 
Die dargestellten Grundlagen bilden die Basis für die Entwicklung und Bewertung der eingesetzten Systems. %todo

\subsection{Der Dobot Magician als Roboterplattform}
\label{sec:dobot_magician-chapter2}
Im folgenden Abschnitt werden die technischen Spezifikationen des Roboterarms vorgestellt und beschrieben. % todo: Diese Informationen sind entscheidend, um die Anforderungen an die im folgenden Kapiteln behandelten Lernansätze zu verstehen und bewerten.
Die Spezifikationen wurden aus dem offiziellen Datenblatt des Dobot Magician entnommen und in Bezug auf die Anforderungen der Montage- und Demontageaufgaben analysiert.

\subsubsection{Technische Spezifikationen und Fähigkeiten}
Der Dobot Magician ist ein kompakter, vierachsiger Roboterarm, der häufig für Bildungszwecke und Forschungsaufgaben eingesetzt wird %todo \cite{}. 
Er verfügt über eine Tischmontage sowie eine integrierte Steuerung in der Basis. Der Arm besteht aus Aluminium und ABS-Kunststoff und besitzt ein Eigengewicht von 3,4\,kg (ohne Werkzeug beziehungsweise Endeffektor) \cite{dobot_datasheet}. Die vier Achsen des Roboterarms sind wie folgt konfiguriert: J1 (Basisrotation), J2 (Schultergelenk), J3 (Ellenbogengelenk) sowie J4 (Handgelenk beziehungsweise Greiferrotation). Diese Konfiguration ermöglicht ein breites Bewegungspektrum, das für Montage- und Demontageaufgaben relevant ist \cite{dobot_datasheet}.

Die maximale Reichweite des Dobot Magician beträgt 320\,mm \cite{dobot_datasheet} und definiert damit den Arbeitsraum des Benchmark-Tasks. 
Die Wiederholgenauigkeit liegt bei $\pm$0,2\,mm \cite{dobot_datasheet} und stellt die Grenze der maximal erreichbaren Genauigkeit sowie der Reproduzierbarkeit des Systems dar. 
%todo: Daraus ergibt sich die Anforderung, dass alle Lernansätze diese Toleranz einhalten oder übertreffen müssen, um innerhalb dieses Bereichs präzise zu agieren und sicherzustellen, dass die Hardware nicht zum limitierenden Faktor wird. 
Die maximale Gelenkgeschwindigkeit beträgt 320°/s bei einer Nutzlast von 250\,g, während das Rotationsservo bis zu 480°/s erreicht \cite{dobot_datasheet}. 

Die maximale Traglast des Roboterarms beträgt 500\,g \cite{dobot_datasheet}, wobei bei einer reduzierten Nutzlast von 250\,g die Höchstgeschwindigkeit und maximale Dynamik erreicht wird. 

Angesteuert wird der Dobot Magician standardmäßig über eine USB-Verbindung. Zur Steuerung wird herstellerseitig die Software 
%todo da DobotLab genutzt wird mitlerweiel: DOBOTStudio eingesetzt, die für Windows und macOS verfügbar ist und grundlegende Programmierung sowie Bewegungssteuerung unterstützt \cite{dobot_datasheet}. 
Die I/O-Schnittstelle an der Basis stellt zwei schaltbare 12\,V- und zwei 5\,V-Ausgänge sowie vier 3,3\,V-Ein-/Ausgänge bereit, von denen zwei als PWM-Ausgänge (Pulsweitenmodulation) und zwei als ADC-Eingänge (Analog-Digital-Converter) konfiguriert sind; zusätzlich sind drei RS232-Schnittstellen verfügbar \cite{dobot_datasheet}. Am Arm selbst stehen zwei schaltbare 12\,V-Ausgänge sowie sieben 3,3\,V-Ein-/Ausgänge zur Verfügung, davon drei PWM- und vier ADC-Kanäle \cite{dobot_datasheet}.

Für die Stromversorgung wird eine Spannung von 100–240\,V\,AC benötigt, die intern auf 12\,V\,DC bei maximal 7\,A gewandelt wird, was zu einer maximalen Leistungsaufnahme von 60\,W führt \cite{dobot_datasheet}. Die zulässige Umgebungstemperatur für den Betrieb liegt zwischen \SI{-10}{\degreeCelsius} und \SI{60}{\degreeCelsius} \cite{dobot_datasheet}.


\subsubsection{Einschränkungen und Grenzen}
\label{sec:dobot-einschraenkungen}
Obwohl der Dobot Magician eine vielseitige Plattform darstellt, ergeben sich durch seine Hardware- und Softwarespezifikationen signifikante Einschränkungen für den Benchmark-Einsatz:

\begin{itemize}
    \item \textbf{Nutzlast und Dynamik:} Die strikte Obergrenze der Traglast von 500\,g begrenzt die Auswahl geeigneter Benchmark-Objekte direkt. Montage- und Demontageobjekte mit einer Masse oberhalb von 500\,g werden daher kategorisch ausgeschlossen. Zudem reduzieren höhere Lasten jenseits der optimalen 250\,g die Dynamik des Systems \cite{dobot_datasheet}.
    \item \textbf{Konnektivität und Software-Integration:} Eine kabellose Kommunikation via WLAN oder Bluetooth ist standardmäßig nicht vorhanden und ausschließlich bei Modellen mit der \enquote{Advanced Communication}-Erweiterung verfügbar \cite{dobot_datasheet}. Auf Softwareebene 
    %todo: bietet die Standardanwendung DOBOTStudio keine native Integration in Simulationsumgebungen oder gängige Middleware-Frameworks wie ROS, was den Aufwand für die Implementierung fortgeschrittener Lernansätze erhöht.
    \item \textbf{Umgebungsbedingungen:} Der Roboterarm ist rein für den Laborbetrieb konzipiert und entspricht lediglich der Schutzklasse IP20. Er ist daher ausschließlich in trockenen Innenräumen zu betreiben \cite{dobot_datasheet} und gegen Staub oder Feuchtigkeit ungeschützt.
\end{itemize}


\subsection{Computervision Grundlagen}
    \subsubsection{Kamera-Modell und Linsenverzerrung}
    \subsubsection{Kamera-Modell und Linsen}
    \subsubsection{Homographie und Pixel-zu-Roboterkoordinaten-Transformation}
    \subsubsection{Kalibrierungsverfahren}

\subsection{Objekterkennung in der Robotik}
    \subsubsection{YOLOv8/HSV für schnelle Vordetektion}
    \subsubsection{Grounding DINO für verfeinerte Erkennung}
    \subsubsection{Coarse-to-Fine Ansatz}

\subsection{Unified Planning für Aufgabenplanung}
    \subsubsection{Grundkonzepte und Funktionsweise}
    \subsubsection{Task-Definition und Optimierung}

\subsection{Mensch-Roboter-Kollaboration}
    \subsubsection{Anforderungen und Sicherheitsaspekte}
    \subsubsection{Schnittstellen und Kommunikation}


\newpage


\section{Methoden}
\subsection{Experimenteller Aufbau}
    \subsubsection{Hardware-Setup (Kamera, Roboter, Workbench, Komponenten)}
    \subsubsection{Kalibrierungsprozess und Validierung}
    \subsubsection{Softwarestack und Implementierung (Python, Libraries, pydobotplus)}

\subsection{Kamera-Verzerrungskorrektur und die Transformation von Pixel- in Roboterkoordinaten}
    \subsubsection{Aufnahme und Verarbeitung der Kalibrierbilder}
    \subsubsection{Berechnung der Distortionsmatrix}
    \subsubsection{Homographie-Berechnung und Koordinatentransformation}
    \subsubsection{Validierung der Kalibriergenauigkeit}

\subsection{Detektion und Klassifizierung von Komponenten}
    \subsubsection{YOLOv8 Vordetektion}
    \subsubsection{Grounding DINO Verfeinerung und Objektklassifizierung}
    \subsubsection{Generierung der Objektlisten mit Positionen und Klassen}

\subsection{Unified Planning zur Generierung von Montage-Sequenzen}
    \subsubsection{Definition der Szenarien [SZENARIO 1 / SZENARIO 2 / SZENARIO 3]}
    \subsubsection{Task-Definition und Task-Labeling (Roboter vs. Mensch)}
    \subsubsection{Plan-Generierung basierend auf erkannten Objekten}
    \subsubsection{Validierung der Machbarkeit}

\subsection{Definition und Umsetzung der Mensch-Aufgaben}
    \subsubsection{Rotation des Fahrzeugs zur Ausrichtung}
    \subsubsection{Werkzeugwechsel und Schraubenarbeiten}
    \subsubsection{Validierung und Fehler-Feedback}
    \subsubsection{Ad-hoc Aufgaben während der Durchführung}

\subsection{Roboter-basierte Greif- und Montage-Operationen}
    \subsubsection{Verfügbare Roboter-Skills und deren Parametrisierung}
    \subsubsection{Greif- und Montagestrategien für verschiedene Komponentenklassen}

 4.6.2.1 Objektklassifizierung und Symmetrie-Analyse
- Klassifizierung der zu montierenden Objekte (Schrauben, Räder, Bumper, Side Covers, Top Addons)
- Symmetrie-Analyse: Welche Objekte sind rotationsinvariant, welche nicht?
- Tabelle: Objekt → Symmetrie → Orientierungsanforderung

 4.6.2.2 2D-Pose Estimation aus Segmentierungsmaske
- Extraktion der Segmentiermaske aus GDINO-Detektionen
- Fitting von Ellipsen zur Bestimmung der Hauptachse
- Berechnung des 2D-Orientierungswinkels (Yaw)
- Confidence-Metriken basierend auf Aspekt-Verhältnis

 4.6.2.3 Strategie-Mapping pro Objektklasse
- Tabelle: Objektklasse → Greif-Strategie → Parameter
- Beispiele:
  * Symmetrische Objekte (Schrauben, Räder): Position-basiert, keine Orientierung nötig
  * Asymmetrische Objekte (Bumper, Side Cover, Top Addon): Position + 2D-Orientierung erforderlich
- Höhen-Offsets für verschiedene Teile

 4.6.2.4 Integration in Manipulations-Pipeline
- Workflow: Detection → Segmentierung → Pose Estimation → Gripper-Ausrichtung → Greifen
- Pseudocode/Algorithmus
- Fehlerbehandlung und Fallback-Strategien

 4.6.2.5 Ergebnisse und Validierung
- Greif-Erfolgsraten pro Objektklasse
- Orientierungs-Genauigkeit (Soll vs. Ist)
- Visualisierungen: Beispiel-Greifposen mit Orientierungsvektoren

    \subsubsection{Kommunikation mit dem Roboter via Python-Bibliothek}
    \subsubsection{Koordination von Roboter- und Mensch-Aufgaben}
    \subsubsection{Error-Handling und Optionen für Mensch-Einspringen}

\subsection{Systemintegration und Gesamtablauf}
    \subsubsection{Datenfluss und Pipeline-Architektur}
    \subsubsection{Timing und Synchronisation zwischen Komponenten}
    \subsubsection{Ablauf eines typischen Montage-Zyklus}


\newpage


\section{Ergebnisse und Evaluierung}

\subsection{Gesamtsystem-Evaluation}
    \subsubsection{[SZENARIO 1]: Erfolgsquote, Durchführungsdauer, kritische Punkte}
    \subsubsection{[SZENARIO 2]: Erfolgsquote, Durchführungsdauer, kritische Punkte}
    \subsubsection{[SZENARIO 3]: Erfolgsquote, Durchführungsdauer, kritische Punkte}
    
    \subsubsection{Robustheit und Flexibilität unter Variation}
        \paragraph{Baseline-Bedingungen (optimale Objektplatzierung)}
        \paragraph{Objekte am Workspace-Rand}
        \paragraph{Veränderte Objektpositionen}
        \paragraph{Kombinierte Variationen}
        \paragraph{Tabellarische Zusammenfassung der Erfolgsquoten}

\subsection{Komponenten-Einzelevaluierung}
    \subsubsection{Detektor-Accuracy: GDINO-Performance pro Komponentenklasse}
    \subsubsection{Kalibrierung: Fehleranalyse bei Pixel-zu-Roboter-Transformation}
    \subsubsection{Roboter-Greifung: Erfolgsquote nach Komponentenklasse und Platzierung}
    \subsubsection{Unified Planning: Plan-Generierungszeit und Task-Komplexität pro Szenario}

\subsection{Analyse und Bottleneck-Identifikation}
    \subsubsection{Hauptprobleme und deren Ursachen}
    \subsubsection{Komponenten-Limitierungen und Auswirkungen auf das Gesamtsystem}
    \subsubsection{Ad-hoc Aufgaben und praktische Erkenntnisse}

\subsection{Fokus auf die Stärken des Systems}
    \subsubsection{Funktionalität und End-to-End-Erfolg}
    \subsubsection{Robustheit unter variierenden Bedingungen}
    \subsubsection{Mehrwert der Mensch-Roboter-Kollaboration}

\newpage

\section{Fazit und Ausblick}
\label{sec:fazit}

\subsection{Zusammenfassung der Ergebnisse}
\subsection{Beantwortung der Forschungsfrage}
\subsection{Erfolgreiche Lösung des Gesamtsystems versus einzelne Komponenten-Schwächen}
\subsection{Grenzen dieser Arbeit}

\subsection{Zukünftige Arbeiten}
    \subsubsection{Strukturiertes Error-Handling und Fallback-Mechanismen}
    \subsubsection{Erweiterung auf weitere Szenarien und Komponenten-Klassen}
    \subsubsection{Optimierung einzelner Komponenten (GDINO-Accuracy, Roboter-Präzision, Planung)}
    \subsubsection{Steigerung der Automatisierung (Reduktion Mensch-Tasks)}


\newpage

\printbibliography[heading=bibintoc,title=Literaturverzeichnis]

\end{document}
